\documentclass[11pt]{article}
\usepackage[a4paper, top=18mm, bottom=18mm, left=20mm, right=20mm]{geometry}
\usepackage[T2A]{fontenc}
\usepackage[utf8]{inputenc}
\usepackage[ukrainian]{babel}
\usepackage{graphicx}
\usepackage[export]{adjustbox}
\usepackage{amsmath}
\usepackage{amssymb}
\graphicspath{{/app/Quiz/Scripts/2023/ProgAll/15BinTree/}{/app/Quiz/Scripts/2020/Mod2020_4/BinTree/}{/app/Quiz/Scripts/2021/Mod2021_3/BinTree/}}
\setlength{\parindent}{0pt}
\setlength{\parskip}{4pt}
\pagestyle{empty}
\begin{document}
%begin
{\large\textbf{Контрольна робота}}

\textbf{Варіант \No\,1}\hfill \textit{ПІБ:}\,\underline{\hspace{60mm}}
\vspace{4mm}

%beginitem
1. 
Скільки 2017-елементних підмножин множини натуральних від 1 до 20172017 містять рівно два числа, що менше 172007
%enditem

%beginitem
2. 

a,b,c=7,8,9
def g(a,b,c=6):
    print(a,b,c,end="")

g(b=1,a=5,c=3)
print(a,b,c)


%enditem

%beginitem
3. 

\begin{center}
	\adjustimage{max size={0.8\linewidth}{0.8\paperheight}}
	{../15BinTree/trees_exp/105.png}
\end{center}
Указати послідовність міток вершин дерева, зображеного на рис., що утвориться в результаті обходу дерева в оберненому порядку. Мітки записувати через пробіл.\par Алгоритм починає роботу з кореня (на рис. це верхня вершина).
%answer:105 оберненому

%enditem

%beginitem
4. %goodpath:Scripts/2018/Mod2018_1/01-good.txt
Відмітити все, що є однією правильною лексемою:
( 1 ) For

( 2 ) a1.3

( 3 ) 5_w

( 4 ) Switch

%enditem

%beginitem
5. 
Записати ВИРАЗ, значенням якого є множина, що складається з квадратних коренівцілих чисел від 16 до 2006 (включно), що діляться на 3.
%enditem

%beginitem
6. Що буде виведено за виконання фрагменту коду?

%code
#include <iostream>

int g(int a, int b){
    int c = 51;
    if (a)
        c = 7;
    else if (b <= 1)
         return 5;
    else 
        c = 3;
    return c;
}

int main(){
    std::cout << g(-3, -3);
    return 0;
}
%code

%enditem

%beginitem
7. Що буде виведено за виконання фрагменту коду?

%code
#include <iostream>
using namespace std;

bool f(int n){
    cout<<"f";
    return n<=2;
}

int main(){
    cout<<(f(-3) or f(3));
    return 0;
}
%code

%enditem

%beginitem
8. Що буде виведено за виконання фрагменту коду?

%code
print('R', end=' ')
j=11
while j>=7:
    j+=-2
print(j)
%code


%enditem

%beginitem
9. Що буде виведено за виконання фрагменту коду?

%code
cout << 10 / 10 / 10 / 10;
%code

%enditem

%beginitem
10. 
try:
    res = 9//9/9*9
    if isinstance(res, bool):
        print(f'bool;{res}')
    elif isinstance(res, int):
        print(f'int;{res}')
    elif isinstance(res, float):
        print(f'float;{res:.2f}')
    else:
        print('???')
except:
    print('Z')

%enditem

%beginitem
11. 
a = 7
b = 5
c = 3

def g(b):
global c    a = 2
    b = 5
    c = 4
    return a + b + c

a = 0
b = 4
c = 6

try:
    print(g(a), a, b, c, sep=';')
except:
    print('ERR', a, b, c, sep=';')

%enditem

%beginitem
12. Що буде виведено за виконання фрагменту коду?

%code
print(7.0 <= "False")
%code

%enditem

%beginitem
13. Що буде виведено за виконання фрагменту коду?
code
__b = 0

class B:
    __b = 1
    
    def __init__(self):
        self.__b = 2
        __b = 3
        B.__b = 4

obj = B()
try:
    print(__b, end=' ')
    print(obj.__b, end=' ')
    print(B.__b, end=' ')
    print(obj.__b, end=' ')
    print(obj._B__b, end=' ')
except:
    print('ERR')
code

%enditem

%beginitem
14. Що буде виведено за виконання фрагменту коду?

%code
h = (0,1,2,3,4,5,6,7,8,9)
print(h[-9::-2])
%code


%enditem

%beginitem
15. 
Рядки журналу через двокрапку містять ім'я користувача (ідентифікатор), час події,тип події, код події, наприклад
\begin{verbatim}
s="username : 2022-05-05 13:26 : ERROR : 404"
\end{verbatim}
у коді 
\begin{verbatim}
res = re.fullmatch('???', s) 
s = ''
code_str = ??? 
\end{verbatim}
чим треба замінити кожен з ???, щоб значенням code\_str був код події? 



%enditem

%beginitem
16. Що буде виведено за виконання фрагменту коду?

%code
print('R', end=' ')
h=('a','b','c','d','e',0,1,2)
print(h[-7])
%code


%enditem

%beginitem
17. Що буде виведено за виконання фрагменту коду?

%code
a,b,c=7,5,3
def g(b):
global c    a=2
    b=5
    c=4
    return a+b+c

a,b,c=0,4,6
print(g(a),a,b,c)
%code

%enditem

%beginitem
18. 
h=(0,1,2,3,4,5,6,7,8,9); print(h[-9::-2])
%enditem

%beginitem
19. %goodpath:Scripts/2019/Mod2019_1/Lex/03-good.txt
Відмітити все, що є коректним оператором С++:
( 1 ) x==y==z

( 2 ) return +

( 3 ) x not is y

( 4 ) return None

%enditem

%beginitem
20. 
x,y,z=1,5,3;  y,z,x,x=y,x,7,8;  print(x,y,z)
%enditem

%beginitem
21. 
def f(n): print("f", end=""); return n<=2
   print(f(-3) or f(3))
%enditem

%beginitem
22. 
j=7
   while j<=11: j+=2
   print(j)
%enditem

%beginitem
23. Що буде виведено за виконання фрагменту коду?
code
b = 0
j = 1

class B:
    j = 2
    
    def __init__(self):
global j        self.j = 3
        j = 4
        B.j = 5

obj = B()
print(b, j, B.j, obj.j, sep=' ')
code

%enditem

%beginitem
24. Що буде виведено за виконання фрагменту коду?

%code
print('R', end=' ')
h = ('a','b','c','d','e',0,1,2)
print(h[-7])
%code


%enditem

%beginitem
25. 
Значенням змінної \kw{a} є цілочисловий масив типу $numpy.ndarray$ розмірів $4{\times}4$. На скільки збільшиться сума елементів масиву \kw{a} після виконання
\begin{verbatim}
a.T[-2] += 2
\end{verbatim}


Якщо код некоректний, то в якості відповіді зазначити ERROR.



%enditem

%beginitem
26. 
$a$ є цілочисловим масивом типу $numpy.ndarray$ розмірів $3{\times}3$. Сума елементів масиву $a$ дорівнює 26. Чому дорівнює сума елементів масиву $a$ після виконання 
\begin{verbatim}
a -= 26
\end{verbatim}


Якщо код некоректний, то в якості відповіді зазначити ERROR.



%enditem

%beginitem
27. 
not9.0<=False or not7.0!=9 or 7>=5
%enditem

%beginitem
28. 
h=['0','1','2','3','4']; h[:-2]=[1,2]; print(h)
%enditem

%beginitem
29. 
Змінні \kw{next} та \kw{prev} вузла списку позначають наступний та попередній за ним вузол відповідно. Починаючи з голови, у список записано послідовні цілі числа від~$-29$ до~$22$. Нехай змінна \kw{p} позначає вузол, що зберігає значення~$-3$.
Яке число зберігається у вузлі \kw{p.next.next.next.next.next} ?

%enditem

%beginitem
30. 
def f(arg, n):
    n = 7
    tmp = arg
    tmp[:-2] = 'xy'
    return len(tmp)>3

try:
    h = ['0','1','2','3','4']
    n = 5
    f(h, n)
    print(n, end=';')
    for el in h:
        print(el, end=';')
except:
    print('Z')

%enditem

%beginitem
31. 
h=['0','1','2','3','4']; h+=[2,3]; print(h)
%enditem

%beginitem
32. Що буде виведено за виконання фрагменту коду?
%code

x,y,z=1,5,3
y,z,x,x=y,x,7,8
print(x,y,z)
%code

%enditem

%beginitem
33. Що буде виведено за виконання фрагменту коду?

%code
print(not9.0<=False or not7.0!=9 or 7>=5)
%code

%enditem

%beginitem
34. 

print('R', end=' ')
h=['0','1','2','3','4']
h[:-2]=[1,2]
print(h)


%enditem

%beginitem
35. 
def f(arg, n):
    n = 7
    tmp = arg
    tmp[:-2] = 'xy'
    return len(tmp)>3

try:
    h = ['0','1','2','3','4']
    n = 5
    f(h, n)
    print(n, end=';')
    for el in h:
        print(el, end=';')
except:
    print('ERR')

%enditem

%beginitem
36. Що буде виведено за виконання фрагменту коду?

%code
print(7.0<="False")
%code

%enditem

%beginitem
37. Що буде виведено за виконання фрагменту коду?

%code
#include <iostream>
using namespace std;

int g(int a, int b){
    int c = 51;
    if (a)
        c = 7;
    else if (b <= 1)
         return 5;
    else 
        c = 3;
    return c;
}

int main(){
    cout << g(-3, -3);
    return 0;
}
%code

%enditem

%beginitem
38. 

\begin{center}
	\adjustimage{max size={0.8\linewidth}{0.8\paperheight}}
	{../BinTree/trees_exp/105.png}
\end{center}
Указати послідовність міток вершин дерева, зображеного на рис., що утвориться в результаті обходу дерева в оберненому порядку. Мітки записувати через пробіл.\par Обхід дерева починається з кореня (на рис. це верхня вершина).
%answer:105 оберненому

%enditem

%beginitem
39. Що буде виведено за виконання фрагменту коду?
%code
#include <iostream>
using namespace std;

int f(int &x, int &y){
    x = 7;
    y-= 5;
    return y;
}

int main(){
    int a = 4, b = 3;
    b = f(b, b);
    cout << a << ":" << b <<':';
    {
        int a = 1, b = 8;
        cout << ((b<=4) || ((a-=2) <= 4)) << ':';
        cout << a << ":" << b << ':';
    }
    {
        inta = 2;
        cout << a << ':';
    }
    cout << a << ":" << b << endl;
    return 0;
}
%code


%enditem

%beginitem
40. 
a,b,c=7,5,3
    def g(b):
      global c      a=2
      b=5
      c=4
      return a+b+c
    a,b,c=0,4,6
    print(g(a),a,b,c)
%enditem

%beginitem
41. 
В умовах попередньої задачі було виконано p->next=p->prev->next;

Чому дорівнює значення p->next->next->next->next->n ?
%enditem

%beginitem
42. Що буде виведено за виконання фрагменту коду?

%code
h=['0','1','2','3']
h+=[2,3]
print(h)
%code


%enditem

%beginitem
43. Що буде виведено за виконання фрагменту коду?

%code
9//9/9*9
%code

%enditem

%beginitem
44. 

print('R', end=' ')
j=11
while j>=7:
    j+=-2
print(j)


%enditem

%beginitem
45. Що буде виведено за виконання фрагменту коду?

%code
print('R', end=' ')
h=['0','1','2','3','4']
h[:-2]=[1,2]
print(h)
%code


%enditem

%beginitem
46. Що буде виведено за виконання фрагменту коду?

%code
a = 7
b = 5
c = 3

def g(b):
global c    a-=2
    b=5
    c=4
    return a+b+c

a, b, c = 0, 4, 6
try:
    print(g(a), a, b, c, sep=';')
except:
    print('Z', a, b, c, sep=';')
%code

%enditem

%beginitem
47. Що буде виведено за виконання фрагменту коду?

%code
#include <iostream>
using namespace std;

int a = 7, b = 5, c = 3;

int g(int &b){
int c;    a -= 2;
    b = 5;
    c = 4;
    return a + b + c;
}

int main(){
    inta = 0;
    intb = 4;
    intc = 6;
    cout << g(a) << ':';
    cout << a << ':' << b << ':' << c;
    return 0; 
}
%code

%enditem

%beginitem
48. 
a,b,c=7,5,3
    def g(b):
      global c      a-=2
      b=5
      c=4
      return a+b+c
    a,b,c=0,4,6
    print(g(a),a,b,c)
%enditem

%beginitem
49. Що буде виведено за виконання фрагменту коду?

%code
try:
    j = 11
    while j>=7:
        j += -2
    print(j, end=';')
except:
    print('Z')

%code

%enditem

%beginitem
50. 

\begin{center}
	\adjustimage{max size={0.8\linewidth}{0.8\paperheight}}
	{../BinTree/trees_exp/105.png}
\end{center}
Указати послідовність міток вершин дерева, зображеного на рис., що утвориться в результаті обходу дерева в оберненому порядку. Мітки записувати через пробіл.\par Алгоритм починає роботу з кореня (на рис. це верхня вершина).
%answer:105 оберненому

%enditem

%beginitem
51. 
try:
    j = 11
    while j>=7:
        j += -2
    print(j, end=';')
except:
    print('Z')

%enditem

%beginitem
52. Що буде виведено за виконання фрагменту коду?

%code
def g(a,b):
    c=51
    if a:
        c=7
    elif b<=1:
         return 5
    else: 
        c=3
    return c

print(g(-3,-3))
%code

%enditem

%beginitem
53. Що буде виведено за виконання фрагменту коду?

%code
#include <iostream>
using namespace std;

int g(int b){
    int v = 51;
    if (b) 
        v = 7;
    else if (b <= 1)
         return 5;
    else
         v = 3;
    return v;
}

int main(){
    cout << g(-3);
    return 0;
}
%code


%enditem

%beginitem
54. %goodpath:Scripts/2018/Mod2018_1/03-good.txt
Відмітити все, що є коректним оператором С++:
( 1 ) int x,y;

( 2 ) if ('x'<>'y') a=6;

( 3 ) double n.m;

( 4 ) {int x,y;}

%enditem

%beginitem
55. Що буде виведено за виконання фрагменту коду?

%code
def g(b):
    v=51
    if b: 
        v=7
    elif b<=1:
         return 5
    else:
         v=3
    return v

print(g(-3))
%code

%enditem

%beginitem
56. Що буде виведено за виконання фрагменту коду?

%code
cout << (!9.0<=false or !7.0!= 9 or 7>=5);
%code

%enditem

%beginitem
57. Що буде виведено за виконання фрагменту коду?

%code
print('R', end=' ')
h = (0,1,2,3,4,5,6,7,8,9)
print(h[-9::-2])
%code


%enditem

%beginitem
58. 

print('R', end=' ')
h=['0','1','2','3','4']
h+=[2,3]
print(h)


%enditem

%beginitem
59. Що буде виведено за виконання фрагменту коду?

%code
print('R', end=' ')
h=['0','1','2','3','4']
h+=[2,3]
print(h)
%code


%enditem

%beginitem
60. Що буде виведено за виконання фрагменту коду?

%code
print(9//9/9*9)
%code

%enditem

%beginitem
61. 
try:
    j = 11
    while j>=7:
        j += -2
    print(j, end=';')
except:
    print('ERR')

%enditem

%beginitem
62. 
try:
    for b in range(-9, -9, -2):
        if b<=4:
            continue
        print(b, end=';');b = -4
        if b<=4:
            break
    else:
        print(b,end=';')
    print(b)
except:
    print('ERR')

%enditem

%beginitem
63. 
7.0<="False"
%enditem

%beginitem
64. Що буде виведено за виконання фрагменту коду?
%Оригинал был утерян. Требует отладки!
%code
__b = 0

class B:
    __b = 1
    
    def __init__(self):
        self.__b = 2
        __b = 3
        B.__b = 4

obj = B()
B.__b = 5
try:
    print(__b, end=' ')
    print(obj.__b, end=' ')
    print(B.__b, end=' ')
    print(obj.__b, end=' ')
    print(obj._B__b)
except:
    print('ERR')
%code

%enditem

%beginitem
65. 
У папці X55 розташовано проект, до якого входить синаксично правильна одиниця трансляції 1.cpp з рядком\par
{\#}include "2.h"
\parФайла 2.h у папці X55 нема, але він є в папці Y7. Що треба зробити, щоб одиниця трансляції 1.cpp, могла успішно відкомпілюватися?Переміщувати, копіювати та змінювати файли з кодом заборонено.\par\vspace{1.5cm}
%enditem

%beginitem
66. 

print('R', end=' ')
for b in range(-9,-9,-2):
    if b<=4:
        continue
        print(b, end=' ')
        b=-4
    if b<=4:
        break
    else:
        print(b, end=' ')
    print(b)

%enditem

%beginitem
67. Що буде виведено за виконання фрагменту коду?

%code
j=7
while j<=11:
    j+=2
print(j, end=' ')
%code


%enditem

%beginitem
68. 
Записати ВИРАЗ, значенням якого є словник, ключами якого є цілі числа від 16 до 2006 (включно), що діляться на 3, а ключам відповідають їх куби 
%enditem

%beginitem
69. 
Написати програму, що запитує у користувача значення дійсних параметрів $x$ та $y$, обчислює для них значення виразу 
$$\frac{\log_{10} x - 43}{(\arctg x - 0.7) (y - 17)}$$ 
та повiдомляє введенi данi та результат обчислення.
 Оцінюється правильність та якість коду.


%enditem

%beginitem
70. %goodpath:Scripts/2019/Mod2019_1/Lex/01-good.txt
Відмітити все, що є однією правильною лексемою:
( 1 ) ;

( 2 ) "a"a"

( 3 ) '\'

( 4 ) +

%enditem

%beginitem
71. 
Змінні $next$ та $prev$ вузла списку позначають наступний та попередній за ним вузол відповідно. Починаючи з голови, у список записано послідовні цілі числа від $-29$ до $22$. Нехай змінна $p$ позначає вузол, що зберігає значення $-3$.
Яке значення зберігається у вузлі $p.next.next.next.next.next$ ?

%enditem

%beginitem
72. 
for b in range(-9,-9,-2):
    if b<=4:
        continue
        print(b,end=' ')
        b=-4
    if b<=4:
        break
    else:
        print(b,end=' ')
    print(b)

%enditem

%beginitem
73. 
(Задача за частиною Пашка А.О. Наводиться в авторській редакції.)\par {\hspace {0.5 cm}}³дбудеться деяка подія $W$ чи ні залежить від двох факторів $A$ та $X$. За результатами статистичного експерименту (100 раз для кожного значення фактора $A$,  200 раз для кожного значення фактора $X$) подія $W$ відбулася:\par {\hspace {0.5 cm}}$(n+17)$ раз для значення $A(a)$, $(n+21)$ раз для значення $A(b)$, $(n+50)$ раз для значення $A(c)$, $(n+57)$ раз для значення $A(d)$, та відповідно $(2n+100)$ раз для значення $X(x)$, $(4n+17)$ раз для значення $X(y)$, $(n+150)$ раз для значення $X(z)$.\par
{\hspace {0.5 cm}}Знайти ймовірність того що подія $W$ не відбулася (умовна ймовірність) при значеннях факторів $A(d)$, $X(x)$. Фактори є незалежними, $n$ – номер цього екзаменаційного білета.\par
%enditem

%beginitem
74. 
for b in range(-9,-9,-2):
    if b<=4:
        continue
        print(b,end=' ')
        b=-4
    if b<=4:
        break
else:
    print(b, end=' ')
print(b, end=' ')

%enditem

%beginitem
75. 
Записати ВИРАЗ, значенням якого є список, що складається з квадратних коренівцілих чисел від 16 до 2006 (включно), що діляться на 3, записаних в оберненому порядку.
%enditem

%beginitem
76. 

def g(b):
    v=51
    if b: 
        v=7
    elif b<=1:
         return 5
    else:
         v=3
    return v

print(g(-3))

%enditem

%beginitem
77. 
Записати ВИРАЗ, значенням якого є словник, ключами якого є цілі числа від 16 до 2006 (включно), що діляться на 3, а ключам відповідають їх куби 
%enditem

%beginitem
78. 
j=11
   while j>=7: j+=-2
   print(j)
%enditem

%beginitem
79. Що буде виведено за виконання фрагменту коду?
%code
for b in range(-9,-9,-2):
    if b<=4:
        continue
        print(b,end=' ')
        b=-4
    if b<=4:
        break
else:
    print(b, end=' ')
print(b, end=' ')
%code

%enditem

%beginitem
80. 
Проект складається з од���ї одиниці трансляції 1.cpp, яка починається таким фрагментом\par
long g(){ return myFunc();}\par
Функція myFunc(); означена в 2.cpp.\par
Що треба зробити для того, щоб 1.cpp успішно відкомпілювався?\par\vspace{1.5cm}
Що треба зробити для того, щоб за проектом зібрався виконуваний код?\par

%enditem

%beginitem
81. Що буде виведено за виконання фрагменту коду?

%code
cout << (!9.0 <= false or !7.0 != 9 or 7 >= 5);
%code

%enditem

%beginitem
82. 
Записати ВИРАЗ, значенням якого є список, що складається з квадратних коренівцілих чисел від 16 до 2006 (включно), що діляться на 3, записаних в оберненому порядку.
%enditem

%beginitem
83. Що буде виведено за виконання фрагменту коду?

%code
for b in range(-4,-4+3,-2):
    if b<=4:
        continue
    print(b, end=' ')
    b=-4
    if b<=4:
        break
else:
    print(b, end=' ')
print(b, end=' ')
%code

%enditem

%beginitem
84. 
try:
    res = 4**-0.5*-1
    if isinstance(res, bool):
        print(f'bool;{res}')
    elif isinstance(res, int):
        print(f'int;{res}')
    elif isinstance(res, float):
        print(f'float;{res:.2f}')
    else:
        print('???')
except:
    print('ERR')

%enditem

%beginitem
85. Що буде виведено за виконання фрагменту коду?

%code
print(9.0isFalseis7.0is9)
%code

%enditem

%beginitem
86. Що буде виведено за виконання фрагменту коду?

%code
4**-0.5*-1
%code

%enditem

%beginitem
87. 
Знайти кількість відношень на n-елементній множині, які неантирефлексивні або неантисиметричні
%enditem

%beginitem
88. Що буде виведено за виконання фрагменту коду?
%code
if (7 <= 7)
    cout << "h";
else
    cout << "h";
%code


%enditem

%beginitem
89. Що буде виведено за виконання фрагменту коду?

%code
7.0<="False"
%code

%enditem

%beginitem
90. 

a,b,c=7,5,3
def g(b):
global c    a-=2
    b=5
    c=4
    return a+b+c

a,b,c=0,4,6
print(g(a),a,b,c)

%enditem

%beginitem
91. Що буде виведено за виконання фрагменту коду?

%code
cout << (9.0 <= false != 7.0 >= 9);
%code

%enditem

%beginitem
92. 
a = 7
b = 5
c = 3

def g(b):
global c    a = 2
    b = 5
    c = 4
    return a + b + c

a = 0
b = 4
c = 6

try:
    print(g(a), a, b, c, sep=';')
except:
    print('Z', a, b, c, sep=';')

%enditem

%beginitem
93. Що буде виведено за виконання фрагменту коду?

%code
print('R', end=' ')
j=7
while j<=11:
    j+=2
print(j)
%code


%enditem

%beginitem
94. 
try:
    res = not9.0<=False or not7.0!=9 or 7>=5
    if isinstance(res, bool):
        print(f'bool;{res}')
    elif isinstance(res, int):
        print(f'int;{res}')
    elif isinstance(res, float):
        print(f'float;{res:.2f}')
    else:
        print('???')
except:
    print('Z')

%enditem

%beginitem
95. 
try:
    res = 7.0j<="False"
    if isinstance(res, bool):
        print(f'bool;{res}')
    elif isinstance(res, int):
        print(f'int;{res}')
    elif isinstance(res, float):
        print(f'float;{res:.2f}')
    else:
        print('???')
except:
    print('ERR')

%enditem

%beginitem
96. 
Написати програму, що запитує у користувача значення дійсних параметрів $x$ та $y$, обчислює для них значення виразу 
$$\frac{\log_{10} x - 43}{(\arctg x - 0.7) (y - 17)}$$ 
та повiдомляє введенi данi та результат обчислення.


%enditem

%beginitem
97. Що буде виведено за виконання фрагменту коду?

%code
cout << 9 / 9 / 9 / 9;
%code

%enditem

%beginitem
98. 
a = 7
b = 5
c = 3

def g(b):
global c    a -= 2
    b = 5
    c = 4
    return a + b + c

a = 0
b = 4
c = 6

try:
    print(g(a), a, b, c, sep=';')
except:
    print('ERR', a, b, c, sep=';')

%enditem

%beginitem
99. Що буде виведено за виконання фрагменту коду?

%code
print('R', end=' ')
h=(0,1,2,3,4,5,6,7,8,9)
print(h[-9::-2])
%code


%enditem

%beginitem
100. 

try:
    print(7,end="")
    print(5//0.0,end="")
    print(3,end="")
except BaseException: print(0,end="")
except KeyboardInterrupt: print(4,end="")
else: print(6,end="")
finally: print(2,end="")

%enditem

%end
\newpage
\end{document}

